% !TEX root = ../main.tex

\chapter{The Interval Type and Simplices}
\label{chapter-interval}

The interval type plays a pivotal role in SDT in two aspects: It illustrates a primitive homotopy type, as its name suggests, a closed interval, which allows us to define higher-dimensional structures like cubes and simplices, and to talk about homotopy theory in this framework. Meanwhile, it classifies certain subobjects, namely, the open sets on which the partial functions are studied in SDT. This chapter describes how these two aspects of the interval type interact on the properties of simplices, which are important when illustrating information order in homotopy.

\section{The Interval Type}

We denote the interval type as $\I$, along with its two ends as $0$ and $1$, respectively. The existence of the interval type is postulated in the theory. To rule out trivial models, we require $0\ne1$. For any $i\in\I$, denote $\sem{i}:=(i=1)$. To make $\I$ a classifier, we expect $\I$ to be a subuniverse of the universe of propositions $\Omega$, which leads to the following axiom \footnote{Axioms about the interval type will always be postulated throughout the thesis unless otherwise indicated.}:
\begin{axiom}
    \label{ax-prop}
    $\sem{-}:\I\to\Omega$ is an embedding that sends $1\mapsto\top$ and $0\mapsto\bot$.
\end{axiom}

We also observe that $(\Omega,\wedge,\vee,\to,\bot,\top)$ is a bounded distributive lattice, where $\to$ serves as the partial order induced by the lattice. We further expect this embedding to be an embedding of lattice, and as a result, the interval type seen as a subuniverse of propositions will be closed under conjunction and disjunction. 
\begin{axiom}
    \label{ax-lattice}
    $\I$ has a bounded distributive lattice $(\I,\sqcap,\sqcup,\sqsubseteq,0,1)$ structure satisfying
    $$\begin{aligned}
        \sem{i\sqcap j}&=\sem{i}\wedge\sem{j} \\
        \sem{i\sqcup j}&=\sem{i}\vee\sem{j} \\
        i\sqsubseteq j&=\sem{i}\to\sem{j} \\
    \end{aligned}$$
\end{axiom}
It should be emphasised that the name of the interval type does not suggest that the partial order defined here is necessarily a linear order.

\section{The Partial Map Classifier}
We now illustrate how the interval type classifies the subsets and, in addition, the partial maps.
\begin{definition}
    For any $A$, the function type $A\to\I$ is called the \textbf{observational algebra}\footnote{The naming of this concept is attributed to \cite{sterling_domains_2025}.} of $A$, denoted $\Op(A)$. Each element of $\chi:\Op(A)$ identifies a subset $U$ of $A$ defined as
    $$U:=\set{a\in A:\chi(a)=1}=\sum_{a:A}\sem{\chi(a)}$$
    Such a subset is called an \textbf{open subset} of $A$. For simplicity, we may use $U:\Op(A)$ or $U\in\Op(A)$ to refer to the subset itself. The observational algebra is made a bounded distributive lattice by taking pointwise joins and meets.
\end{definition}
One may immediately see that the lattice of the interval induces a bounded distributive lattice on any observational algebra by pointwise operations, which further guarantees that the open sets are closed under finite unions and intersections and that the empty and universal sets are always open. This observation is similar to the definition of the frame of open sets in topology, despite the fact that we so far know nothing about whether open sets are still closed under arbitrary unions. Nevertheless, we still see the observational algebra $\Op(A)$ as the topology of $A$.

With the open set classifier, we can now classify the partial maps defined in the open sets, which is to find a functor $L$ such that $A\to LB$ gives exactly the type of partial maps from $A$ to $B$.

\begin{definition}
    The \textbf{lifting functor} $L$ is defined as
    $$LA:=\sum_{i:\I}(\sem{i}\to A)$$
    that acts on the function $f:A\to B$
    $$Lf:=\lambda(i,x).(i,\lambda\phi.f(x(\phi))):LA\to LB$$
\end{definition}
\begin{theorem}
    \label{thm-classifier}
    The lifting functor classifies partial maps in open subsets. Specifically, for any $A$ and $B$, the set of all partial functions from $A$ to $B$ is isomorphic to $A\to LB$.
\end{theorem}
\begin{proof}
    $$\begin{aligned}
        \sum_{U:\Op(A)}(U\to B)&=\sum_{\chi:\Op(A)}\sum_{a:A}\sem{\chi(a)}\to B \\
        &\cong\sum_{\chi:\Op(A)}\prod_{a:A}\sem{\chi(a)}\to B \\
        &\cong\prod_{a:A}\sum_{\chi:A\to\I}\sem{\chi(a)}\to B \\
        &\cong\prod_{a:A}\sum_{i:\I}\sem{i}\to B \\
        &=A\to LB
    \end{aligned}$$
    Hence, $A\to LB$ classifies the partial maps between $A$ and $B$.
\end{proof}
The lifting functor can be seen as a synthetic counterpart to the maybe monad in the sense that every $LA$ has an element $\bot_A:=(0,\elim_\bot)$ similar to \verb|Nothing| and an embedding
$$\eta_A:A\hookrightarrow LA:=\lambda x.(1,\lambda\_.x)$$
similar to \verb|Just|. In fact, one may consider an model of SDT where $\I$ is just the Boolean type, and the lifting functor in this model is exactly the maybe monad.

\section{The Simplices}
\begin{definition}
    The \textbf{$n$-simplex} $\Delta^n$ in SDT is defined as
    $$\Delta^n:=\set{(i_0,\dots,i_{n-1})\in\I^n:i_0\sqsupseteq\dots\sqsupseteq i_{n-1}}$$
    Specially, $\Delta^0:=\top$ and $\Delta^{-1}:=\bot$.
\end{definition}
Sorting the interval elements in descending order seems counterintuitive, since doing the other way is more common and will illustrate exactly the same structure. The reason we do so will be revealed later, but we still use $\Delta_n$ to denote the $n$-simplex with reversed permutation (i.e. in ascending order) of the interval elements.

One surprising result in SDT shows the connection between the simplices and the lifting functor. One can notice this connection from some simple observations.
\begin{lemma}
    \label{lemma-l-delta-0}
    $L\bot\cong\top$.
\end{lemma}
\begin{proof}
    For any $(i,x)\in L\bot$, $x:\sem{i}\to\bot=\sem{i}\to\sem{0}$, which means $i\sqsubseteq0$. Meanwhile, $0$ is the minimum in $\I$, so we must have $i=0$, and therefore $(0,\rec_\bot)$ is the only element in $L\bot$.
\end{proof}
\begin{lemma}
    \label{lemma-l-delta-1}
    $L\top\cong\I$.
\end{lemma}
\begin{proof}
    Denote the only element of $\top$ by $*$. The isomorphism is given as
    $$\begin{aligned}
        (i,x)&\mapsto i \\
        i&\mapsto(i,\lambda\_.*)
    \end{aligned}$$
    which can be easily verified.
\end{proof}
The observations above give rise to the conjecture that $\Delta^n\cong L^n\top\cong L^{n+1}\bot$. However, difficulties arise when we try to show, for example, $L\I\cong\Delta^2$: We have no measures to construct a tuple of two interval elements from an instance of $(i,x)\in L\I$ since we cannot apply $x$ without knowing anything about $i$. It turns out that we need one more axiom to reach the proof.

\begin{definition}
    \label{def-sigma}
    For any type $A$, an \textbf{$L$-algebra} on $A$ is a function $LA\to A$, and a \textbf{$L$-coalgebra} on $A$ is a function $A\to LA$.
\end{definition}
\begin{axiom}
    \label{ax-l-alg}
    There exists a $L$-algebra $\sigma$ on $\I$ satisfying
    $$\sem{\sigma(i,j)}=\sum_{\phi:\sem{i}}\sem{j(\phi)}$$
    which is called the \textbf{$\sigma$-structure} on $\I$.
\end{axiom}
This axiom suggests that $\I$ as a subuniverse of propositions is closed under $\Sigma$-types. A type like $\I$ that inherit the algebra of propositions via an embedding is called a dominance \cite{reus_general_1997}. It is often additionally required that $\sem{i}\leftrightarrow\sem{j}\implies i=j$, yet this is redundant if we assume $\Omega$ to be the universe of h-props. The intuitive interpretation of the $\sigma$-structure leads to some useful lemmata.
\begin{lemma}
    \label{lemma-sigma-le}
    For any $(i,j)\in L\I$, $\sigma(i,j)\sqsubseteq i$.
\end{lemma}
\begin{proof}
    By Axiom \ref{ax-l-alg}, the left projection $\pi_0$ of $\sum_{\phi:\sem{i}}\sem{j(\phi)}$ has the type $$\pi_0:\sem{\sigma(i,j)}\to\sem{i}$$
    Thus by Axiom \ref{ax-lattice}, $\sigma(i,j)\sqsubseteq i$.
\end{proof}
\begin{lemma}
    \label{lemma-sigma-meet}
    For any $i,j\in\I$, $\sigma(i,\lambda\_.j)=i\sqcap j$.
\end{lemma}
\begin{proof}
    By Axiom \ref{ax-lattice} and \ref{ax-l-alg}, we have
    $$\sem{\sigma(i,\lambda\_.j)}=\sum_{\phi:\sem{i}}\sem{j}=\sem{i}\wedge\sem{j}=\sem{i\sqcap j}$$
    Since $\sem{-}$ is monic by Axiom \ref{ax-prop}, $\sigma(i,\lambda\_.j)=i\sqcap j$.
\end{proof}
\begin{lemma}
    \label{lemma-sigma-eta}
    For any $(i,j)\in L\I$ and $\phi\in\sem{i}$, $j(\phi)=\sigma(i,j)$.
\end{lemma}
\begin{proof}
    Whenever the proposition $\sem{i}\in\Omega$ is inhabited, we can always replace $i$ with $1$ and any $\phi\in\sem{i}$ with $*\in\top$. Thus, we have
    $$\sigma(i,j)=\sigma(1,j)=\sigma(1,\lambda\_.j(*))=1\sqcap j(*)=j(*)=j(\phi)$$
\end{proof}

Although the $\sigma$-structure is powerful enough to prove $L\I\cong\Delta^2$, we need to generalise it to higher dimensions to prove our conjecture in any dimension.
\begin{definition}
    \label{def-sigma-n}
    The $\sigma$-structure on $\I$ can be generalised to $L$-algebras on any $n$-cube $\I^n$ given by
    $$\sigma_n(i,j:\sem{i}\to\I^n):=(\sigma(i,\pi_0\circ j),\dots,\sigma(i,\pi_{n-1}\circ j))$$
    which is called the \textbf{$\sigma$-structure} on $\I^n$.
\end{definition}
Likewise, we have the higher-dimensional counterparts for Lemma \ref{lemma-sigma-le}, \ref{lemma-sigma-meet}, and \ref{lemma-sigma-eta}.
\begin{lemma}
    \label{lemma-sigma-le-n}
    For any $(i,j)\in L\I^n$ and any $0\le k< n$, we have
    $$\pi_k(\sigma_n(i,j))\sqsubseteq i$$
\end{lemma}
\begin{lemma}
    \label{lemma-sigma-meet-n}
    For any $i\in\I$ and $j=(j_0,\dots,j_{n-1})\in\I^n$, we have
    $$\sigma(i,\lambda\_.j)=(i\sqcap j_0,\dots,i\sqcap j_{n-1})$$
\end{lemma}
\begin{lemma}
    \label{lemma-sigma-eta-n}
    For any $(i,j)\in L\I^n$ and $\phi\in\sem{i}$, $j(\phi)=\sigma_n(i,j)$.
\end{lemma}
All the above lemmata are direct consequences of Definition \ref{def-sigma-n} and the corresponding $1$-dimensional cases.

The $\sigma$-structures above are defined on cubes. Before using them to prove our conjecture, we need to restrict it onto simplices.
\begin{lemma}
    The $\sigma$-structure on $\I^n$ defines a $L$-algebra on $\Delta^n\subseteq\I^n$ by restriction.
\end{lemma}
\begin{proof}
    Let $(i,j)\in L\Delta^n$ where $n\ge2$. For any $0\le k\le n-2$,
    $$\forall(\phi:\sem{i}).\pi_k(j(\phi))\sqsupseteq\pi_{k+1}(j(\phi))$$
    Meanwhile, by Definition $\ref{def-sigma}$ and $\ref{def-sigma-n}$, we have
    $$\sem{\pi_k(\sigma_n(i,j))}=\sem{\sigma_n(i,\pi_k\circ j)}=\sum_{\phi:\sem{i}}\sem{\pi_k(j(\phi))}$$
    and similarly
    $$\sem{\pi_{k+1}(\sigma_n(i,j))}=\sum_{\phi:\sem{i}}\sem{\pi_{k+1}(j(\phi))}$$
    Axiom \ref{ax-lattice} assures that the function type $\sem{\pi_{k+1}(j(\phi))}\to\sem{\pi_k(j(\phi))}$ is inhabited, and thus so is $\sem{\pi_{k+1}(\sigma_n(i,j))}\to\sem{\pi_k(\sigma_n(i,j))}$. Therefore, $\sigma_n(i,j)$ gives a descending sequence and is an element of $\Delta^n$.
\end{proof}

We can now prove the conjecture that $\Delta^n\cong L^n\top$
\begin{lemma}[\texttt{SemiLattice.L□↓≡□↓}]
    \label{lemma-l-delta-n}
    For any $n\ge1$, $L\Delta^n\cong\Delta^{n+1}$.
\end{lemma}
\begin{proof}
    We give the isomorphism as follows:
    $$\begin{aligned}
        f(i,j)&:=(i,\sigma_n(i,j))&:L\Delta^n\to\Delta^{n+1} \\
        g(i)&:=(i_0,\lambda\_.(i_1,\dots,i_n))&:\Delta^{n+1}\to L\Delta^n \\
    \end{aligned}$$
    By Lemma $\ref{lemma-sigma-le-n}$, the concatenation $(i,\sigma(i,j))$ is descending and is an element of $\Delta^{n+1}$. This ensures that $f$ is well defined.

    We first show $f\circ g=\id$. Take any $(i_0\sqsupseteq\dots\sqsupseteq i_n)\in\Delta^{n+1}$. We can see
    $$\begin{aligned}
        f(g(i_0,\dots,i_n))=&(i_0,\sigma_n(i_0,\lambda\_.(i_1,\dots,i_n))) \\
        =&(i_0,i_0\sqcap i_1,\dots,i_0\sqcap i_n) \\
        =&(i_0,i_1,\dots,i_n) \\
    \end{aligned}$$
    where $i_0\sqcap i_k=i_k$ for any $0\le k\le n$ since $i_0\sqsupseteq i_k$.

    We then show $g\circ f=\id$. Take any $(i,j)\in L\Delta^n$. We can see
    $$g(f(i,j))=(i,\lambda\_.\sigma_n(i,j))$$
    We will show that $j=\lambda\_.\sigma_n(i,j)$. By the extensionality of functions, this is equivalent to showing that $j(\phi)=\sigma_n(i,j)$ for any $\phi\in\sem{i}$, which is exactly Lemma \ref{lemma-sigma-eta-n}.
\end{proof}
Combining Lemma \ref{lemma-l-delta-0}, \ref{lemma-l-delta-1}, and \ref{lemma-l-delta-n} together, we can immediately reach our goal.
\begin{theorem}[\texttt{SemiLattice.Δ≡□↓}]
    \label{thm-l-delta}
    By iteratively composing the maps defined in Lemma \ref{lemma-l-delta-n}, we can get the following isomorphism:
    $$\Delta^n\cong L^n\top\cong L^{n+1}\bot$$
\end{theorem}
This theorem relates the simplices in homotopy theory to the lifting functor in domain theory. It will enable us to give a homotopy interpretation of the initial (final) $L$-(co)algebra in the future. Before that, we will first investigate the topology of simplices using the Phoa principle.
